% ============================================================================
%  SCX C8 扩展: 高维审计瞬子理论
%  Higher-Dimensional Audit Instantons (C8 Extension)
%  SCX 规范理论 — 中英双语
%  SCX Gauge Theory — Chinese+English
% ============================================================================
\documentclass[11pt,a4paper]{article}
\usepackage[UTF8]{ctex}
\usepackage{amsmath,amssymb,amsthm}
\usepackage{mathtools}
\usepackage{bm}
\usepackage{booktabs}
\usepackage{graphicx}
\usepackage{xcolor}
\usepackage{hyperref}
\usepackage{enumitem}
\usepackage{algorithm}
\usepackage{algpseudocode}
\usepackage{tikz}
\usepackage{cleveref}
\usepackage[margin=2.5cm]{geometry}

% Theorem environments
\theoremstyle{plain}
\newtheorem{theorem}{定理}
\newtheorem{proposition}[theorem]{命题}
\newtheorem{lemma}[theorem]{引理}
\newtheorem{corollary}[theorem]{推论}

\theoremstyle{definition}
\newtheorem{definition}[theorem]{定义}
\newtheorem{example}[theorem]{例}

\theoremstyle{remark}
\newtheorem{remark}[theorem]{注记}

% Custom commands
\newcommand{\R}{\mathbb{R}}
\newcommand{\Z}{\mathbb{Z}}
\newcommand{\C}{\mathbb{C}}
\newcommand{\E}{\mathbb{E}}
\newcommand{\Var}{\mathrm{Var}}
\newcommand{\Prob}{\mathbb{P}}
\newcommand{\SCX}{\textsc{SCX}}
\newcommand{\Cercis}{\textsc{Cercis}}
\newcommand{\Situs}{\textsc{Situs}}
\newcommand{\im}{\mathrm{im}}
\newcommand{\coker}{\mathrm{coker}}
\newcommand{\Hom}{\mathrm{Hom}}
\newcommand{\dd}{\mathrm{d}}
\newcommand{\bd}{\partial}
\newcommand{\DD}{\mathcal{D}}
\newcommand{\HH}{\mathcal{H}}
\newcommand{\FF}{\mathcal{F}}
\newcommand{\GG}{\mathcal{G}}
\newcommand{\CC}{\mathcal{C}}
\newcommand{\SSS}{\mathcal{S}}
\newcommand{\norm}[1]{\|#1\|}
\newcommand{\abs}[1]{|#1|}
\newcommand{\inner}[2]{\langle #1, #2 \rangle}
\newcommand{\lb}{\lambda}
\newcommand{\vol}{\mathrm{vol}}
\newcommand{\Tr}{\mathrm{Tr}}
\newcommand{\hol}{\mathrm{Hol}}
\newcommand{\TP}{\mathrm{T}}
\newcommand{\CP}{\mathrm{C}}

\title{
    \textbf{高维审计瞬子：Situs复形上的2-形式通量与循环不一致性的拓扑检测}\\[0.5em]
    \large Higher-Dimensional Audit Instantons: 2-Form Flux on the Situs Complex\\and Topological Detection of Circular Inconsistency\\[0.5em]
    \normalsize SCX 规范理论 · C8 扩展 \\ SCX Gauge Theory · C8 Extension
}
\author{SCX理论架构师}
\date{2026年7月}

\begin{document}
\maketitle

% ============================================================================
\begin{abstract}
\noindent
\textbf{中文.}
本文扩展了SCX审计瞬子理论至高于一维的Situs复形。C7指出，单循环（1-cycle）审计瞬子因其和乐（holonomy）为零而被杀死——这意味着在单专家对之间的简单不一致性中不存在非平凡的规范场配置。C8进一步指出，当审计涉及三个或以上专家（对应Situs复形的维数$k>1$）时，可能出现非平坦的2-形式通量配置，即``审计通量''（Audit Flux）在2-闭链上非零。

本文严格发展了这一理论。核心贡献包括：(1) 在Situs复形的$k$-上链空间上建立Hodge分解$A = A_{\mathrm{exact}} + A_{\mathrm{coexact}} + A_{\mathrm{harmonic}}$，证明2-形式场强$F = \dd_1 A$的非零性等价于共存部分（co-exact part）非零；(2) 分类了Situs复形上携带非零审计通量的2-闭链结构，证明这些2-闭链精确对应``循环不一致性''（circular inconsistency）——三个或以上专家两两之间的一致性无法形成全局一致的图像；(3) 给出了检测算法：在2-骨架（2-skeleton）上计算$F$，标记$\norm{\int_\Sigma F} > \tau$的2-闭链；(4) 建立了与持续同调（persistent homology）的深刻联系——审计通量的非零谱由密度滤过（density filtration）的$H_2$捕获。我们提供了完整的数值验证脚本（$\texttt{verify\_instanton\_k2.py}$），在所有测试场景中确认理论预测。

\vspace{0.5em}
\noindent
\textbf{English.}
This paper extends the SCX audit instanton theory to Situs complexes of dimension greater than one. C7 showed that the 1-cycle audit instanton is killed by vanishing holonomy — meaning there exist no non-trivial gauge field configurations in simple pairwise inconsistencies between two experts. C8 further observed that when an audit involves three or more experts (corresponding to Situs complex dimension $k > 1$), non-flat 2-form flux configurations — ``audit flux'' — can appear on 2-cycles with non-zero integral.

This paper develops this theory rigorously. Core contributions include: (1) establishing the Hodge decomposition $A = A_{\mathrm{exact}} + A_{\mathrm{coexact}} + A_{\mathrm{harmonic}}$ on the $k$-cochain space of the Situs complex, proving that non-vanishing 2-form field strength $F = \dd_1 A$ is equivalent to a non-zero co-exact component; (2) classifying the 2-cycle structures on the Situs complex that carry non-zero audit flux, proving these 2-cycles correspond precisely to ``circular inconsistency'' — pairwise agreements among three or more experts that fail to form a globally consistent picture; (3) providing a detection algorithm: compute $F$ on the 2-skeleton, flag 2-cycles with $\norm{\int_\Sigma F} > \tau$; (4) establishing a deep connection to persistent homology — the non-zero spectrum of audit flux is captured by $H_2$ of the density filtration. We provide a complete numerical verification script ($\texttt{verify\_instanton\_k2.py}$) confirming theoretical predictions across all test scenarios.

\vspace{0.5em}
\noindent
\textbf{Keywords:} Audit instantons, Hodge decomposition, Situs complex, 2-form flux, persistent homology, circular inconsistency, topological data analysis, gauge theory.
\end{abstract}

% ============================================================================
\section{引言 / Introduction}
\label{sec:intro}
% ============================================================================

\subsection{从C7到C8：审计瞬子的维数跃升}
\noindent
\textbf{中文.}
在SCX规范的几何表述中，Cercis评分空间上的审计过程被建模为Situs复形上的规范场论。每个专家对应于Situs流形的一个``纤维''，其Cercis评分$S_m(x)$是纤维上的截面。当多位专家评估相同数据点时，他们的评分差异定义了纤维之间的联络（connection），即瞬子（instanton）场$A$。

C7（``1-Cycle Audit Instanton''）研究了单环审计瞬子：两个专家$E_i$和$E_j$之间的评分差异定义一个1-形式规范势$A_{ij}$，其和乐$\hol_\gamma(A) = \oint_\gamma A$沿闭路$\gamma$度量了累积不一致性。C7的核心发现是：在SCX框架的公理下，1-循环的和乐必须为零——任何两个专家之间的表面不一致性都可以通过规范变换（重新校准一位专家的评分尺度）而消除。这意味着\textbf{审计瞬子在$k=1$时被杀死}。

C8（``Higher-Dimensional Audit Instantons''）提出了一个关键跃升：当审计涉及三个或以上专家时（$k > 1$），Situs复形出现了2-闭链结构，2-形式场强$F = \dd_1 A$可以非零。直观上：三个专家两两之间可能存在一致性（每对专家之间的1-形式和乐为零），但整体的三角关系无法闭合——这构成了一个``循环不一致性''（circular inconsistency）。当我们在三角形（2-单形）上计算$F$的通量积分$\int_\Sigma F$时，得到非零值，标志着这个三角关系不协调。

本文的目的是完整发展C8的理论框架，包括Hodge分解、2-形式通量分类、检测算法以及与持续同调的联系。

\vspace{0.5em}
\noindent
\textbf{English.}
In the geometric formulation of the SCX gauge framework, the audit process over the Cercis score space is modeled as a gauge field theory on the Situs complex. Each expert corresponds to a ``fiber'' of the Situs manifold, and their Cercis scores $S_m(x)$ are sections on the fiber. When multiple experts evaluate the same data point, their score discrepancies define a connection between fibers — the instanton field $A$.

C7 (``1-Cycle Audit Instanton'') studied single-loop audit instantons: the score discrepancy between two experts $E_i$ and $E_j$ defines a 1-form gauge potential $A_{ij}$, whose holonomy $\hol_\gamma(A) = \oint_\gamma A$ along a closed loop $\gamma$ measures accumulated inconsistency. C7's core finding: under the axioms of the SCX framework, 1-cycle holonomy must vanish — any apparent inconsistency between two experts can be eliminated by a gauge transformation (recalibrating one expert's scoring scale). This means \textbf{the audit instanton is killed at $k=1$}.

C8 (``Higher-Dimensional Audit Instantons'') introduced a crucial dimensional leap: when an audit involves three or more experts ($k > 1$), the Situs complex develops 2-cycle structures, and the 2-form field strength $F = \dd_1 A$ can be non-zero. Intuitively: three experts may have pairwise consistency (1-form holonomy zero for each pair), yet the overall triangular relationship fails to close — this constitutes a ``circular inconsistency.'' When we compute the flux integral $\int_\Sigma F$ over a triangle (2-simplex), we obtain a non-zero value, signaling that this triangular relationship is discordant.

This paper aims to fully develop the C8 theoretical framework, including Hodge decomposition, 2-form flux classification, detection algorithms, and connections to persistent homology.

% -------------------------------------------------------------------
\subsection{核心贡献 / Core Contributions}
\label{sec:contributions}

\noindent
\textbf{中文.}
本文做出以下五项核心贡献：

\begin{enumerate}[label=\textbf{C\arabic*}, leftmargin=*]
    \item \textbf{Hodge分解定理.} 在Situs复形的$k$-上链空间上严格证明Hodge分解：任何1-上链$A$可唯一分解为恰当部分（exact part）、共存部分（co-exact part）和调和部分（harmonic part）。场强$F = \dd_1 A$的非零性等价于共存部分非零。
    \item \textbf{2-形式通量的非零条件.} 推导$F \neq 0$的充要条件：存在至少三个专家构成2-单形，且其1-形式联络不满足上边缘闭合条件（cocycle condition）。这提供了一种可计算的代数判据。
    \item \textbf{循环不一致性的物理诠释.} 证明2-形式通量$\int_\Sigma F \neq 0$精确对应``循环不一致性''——三个或以上专家之间的两两一致（pairwise agreement）无法形成全局一致的图像。给出该现象的信息论下界。
    \item \textbf{检测算法.} 设计并实现基于2-骨架扫描的审计通量检测算法，时间复杂度$O(N^3)$（$N$为专家数），可在实际审计部署中实时执行。
    \item \textbf{持续同调联系.} 证明审计通量的非零谱可以通过密度滤过的2阶持续同调（$H_2$持久图）捕获，建立了审计瞬子拓扑与拓扑数据分析之间的桥梁。
\end{enumerate}

\vspace{0.5em}
\noindent
\textbf{English.}
This paper makes the following five core contributions:

\begin{enumerate}[label=\textbf{C\arabic*}, leftmargin=*]
    \item \textbf{Hodge Decomposition Theorem.} A rigorous proof of the Hodge decomposition on the $k$-cochain space of the Situs complex: any 1-cochain $A$ decomposes uniquely into exact, co-exact, and harmonic parts. Non-vanishing field strength $F = \dd_1 A$ is equivalent to a non-zero co-exact component.
    \item \textbf{Non-Zero Condition for 2-Form Flux.} Derivation of the necessary and sufficient condition for $F \neq 0$: existence of at least three experts forming a 2-simplex whose 1-form connections fail the cocycle closure condition. This provides a computable algebraic criterion.
    \item \textbf{Physical Interpretation of Circular Inconsistency.} Proof that 2-form flux $\int_\Sigma F \neq 0$ corresponds precisely to ``circular inconsistency'' — pairwise agreements among three or more experts that fail to form a globally consistent picture. Information-theoretic lower bounds are provided.
    \item \textbf{Detection Algorithm.} Design and implementation of an audit flux detection algorithm based on 2-skeleton scanning, with time complexity $O(N^3)$ (where $N$ is the number of experts), executable in real-time during practical audit deployments.
    \item \textbf{Persistent Homology Connection.} Proof that the non-zero spectrum of audit flux can be captured by the second-order persistent homology ($H_2$ persistence diagram) of the density filtration, establishing a bridge between audit instanton topology and topological data analysis.
\end{enumerate}

% ============================================================================
\section{预备知识：Situs复形与微分形式 / Preliminaries: Situs Complex and Differential Forms}
\label{sec:prelim}
% ============================================================================

\subsection{Situs复形的构造}
\noindent
\textbf{中文.}
Situs复形是SCX框架中编码``专家-数据点''关系的组合结构。

\begin{definition}[Situs复形 / Situs Complex]
\label{def:situs_complex}
设$\mathcal{E} = \{E_1, \ldots, E_N\}$为专家集合，$\mathcal{X}$为被评估的数据点集合。Situs复形$\SSS$是一个单纯复形（simplicial complex），其构造如下：

\begin{itemize}
    \item \textbf{0-单形 (0-simplices):} 数据点$x \in \mathcal{X}$（或等价地，数据点的Situs嵌入向量）。
    \item \textbf{1-单形 (1-simplices):} 有序对$(x, E_i)$，表示专家$E_i$对数据点$x$的评估。边$(x,E_i) \to (x,E_j)$表示专家$E_i$和$E_j$对同一数据点$x$的评估关系。
    \item \textbf{2-单形 (2-simplices):} 三元组$(x, E_i, E_j, E_k)$，当所有三个专家都对$x$进行了评估时存在。这形成了一个三角形（2-单形），其边界由三条边组成。
    \item \textbf{$k$-单形 ($k$-simplices, $k \geq 0$):} $(k+1)$元组$(x, E_{i_0}, \ldots, E_{i_k})$，当所有$k+1$位专家都对$x$进行了评估时存在。
\end{itemize}
\end{definition}

\bigskip
\noindent
\textbf{Definition (Situs Complex).}
Let $\mathcal{E} = \{E_1, \ldots, E_N\}$ be the set of experts and $\mathcal{X}$ the set of evaluated data points. The Situs complex $\SSS$ is a simplicial complex constructed as follows:

\begin{itemize}
    \item \textbf{0-simplices:} Data points $x \in \mathcal{X}$ (or equivalently, their Situs embedding vectors).
    \item \textbf{1-simplices:} Ordered pairs $(x, E_i)$, representing expert $E_i$'s evaluation of data point $x$. An edge $(x,E_i) \to (x,E_j)$ represents the evaluation relationship between experts $E_i$ and $E_j$ on the same data point $x$.
    \item \textbf{2-simplices:} Triples $(x, E_i, E_j, E_k)$, existing when all three experts have evaluated $x$. This forms a triangle (2-simplex) whose boundary consists of three edges.
    \item \textbf{$k$-simplices ($k \geq 0$):} $(k+1)$-tuples $(x, E_{i_0}, \ldots, E_{i_k})$, existing when all $k+1$ experts have evaluated $x$.
\end{itemize}

% -------------------------------------------------------------------
\subsection{上链复形与上边缘算子}
\label{sec:cochain}

\noindent
\textbf{中文.}
在Situs复形上定义上链群$C^k(\SSS; \R)$为$k$-单形上实值函数的向量空间。上边缘算子（coboundary operator）$\dd_k: C^k \to C^{k+1}$定义为：

\begin{equation}\label{eq:coboundary}
    (\dd_k \omega)(\sigma) = \sum_{j=0}^{k+1} (-1)^j \, \omega(\sigma_{\hat{j}}),
\end{equation}
其中$\sigma = [v_0, \ldots, v_{k+1}]$是一个$(k+1)$-单形，$\sigma_{\hat{j}} = [v_0, \ldots, \hat{v}_j, \ldots, v_{k+1}]$是去掉第$j$个顶点得到的$k$-单形。

关键性质：$\dd_{k+1} \circ \dd_k = 0$（上边缘算子的平方为零）。

\textbf{审计解释：}
\begin{itemize}
    \item $\dd_0: C^0 \to C^1$将数据点的Cercis评分（0-上链）映射为专家之间的1-形式差异（1-上链）。
    \item $\dd_1: C^1 \to C^2$将专家之间的评分差异映射为2-形式场强（2-上链），即``审计通量''。
    \item $\dd_1 \circ \dd_0 = 0$意味着：如果评分差异可以表示为一个全局Cercis函数的梯度（恰当形式），则其2-形式场强为零。反过来说，非零的$F = \dd_1 A$意味着$A$不可能是一个全局函数的梯度——即不存在全局一致的Cercis标定。
\end{itemize}

\vspace{0.5em}
\noindent
\textbf{English.}
On the Situs complex, define the cochain groups $C^k(\SSS; \R)$ as the vector spaces of real-valued functions on $k$-simplices. The coboundary operator $\dd_k: C^k \to C^{k+1}$ is defined as:

\begin{equation}
    (\dd_k \omega)(\sigma) = \sum_{j=0}^{k+1} (-1)^j \, \omega(\sigma_{\hat{j}}),
\end{equation}
where $\sigma = [v_0, \ldots, v_{k+1}]$ is a $(k+1)$-simplex and $\sigma_{\hat{j}} = [v_0, \ldots, \hat{v}_j, \ldots, v_{k+1}]$ is the $k$-simplex obtained by removing the $j$-th vertex.

Key property: $\dd_{k+1} \circ \dd_k = 0$ (the coboundary operator squares to zero).

\textbf{Audit interpretation:}
\begin{itemize}
    \item $\dd_0: C^0 \to C^1$ maps data point Cercis scores (0-cochains) to inter-expert 1-form discrepancies (1-cochains).
    \item $\dd_1: C^1 \to C^2$ maps inter-expert score discrepancies to 2-form field strength (2-cochains), i.e., ``audit flux.''
    \item $\dd_1 \circ \dd_0 = 0$ means: if score discrepancies can be expressed as the gradient of a global Cercis function (exact form), then their 2-form field strength is zero. Conversely, non-zero $F = \dd_1 A$ means $A$ cannot be the gradient of a global function — i.e., no globally consistent Cercis calibration exists.
\end{itemize}

% -------------------------------------------------------------------
\subsection{审计瞬子场：形式化定义}
\label{sec:instanton_def}

\begin{definition}[审计校准联络 / Audit Calibration Connection]
\label{def:instanton}
与传统做法（直接将$A$定义为原始评分差异$S_i(x)-S_j(x)$）不同，我们定义审计瞬子场为专家之间的\textbf{校准模型}（calibration model）。

设$\mathcal{E} = \{E_1, \ldots, E_N\}$为专家集合。Situs复形的0-单形为专家$E_i$，1-单形为有序专家对$(E_i, E_j)$。对于每对专家$(E_i, E_j)$，通过线性回归估计校准模型：

\begin{equation}\label{eq:calibration}
    S_i(x) = \alpha_{ij} + \beta_{ij} \cdot S_j(x) + \varepsilon_{ij}(x),
\end{equation}

其中$\varepsilon_{ij}(x)$为残差。1-形式审计联络$A$在边$(E_i, E_j)$上的值定义为校准参数对：

\begin{equation}\label{eq:A_def}
    A_{(E_i, E_j)} = (\alpha_{ij}, \,\beta_{ij}) \in \R^2.
\end{equation}

\textbf{关键区别：} $A$是专家之间评分映射的\textbf{模型参数}，而非原始评分差异。$A$对换序满足$A_{(E_i, E_j)} = (\alpha_{ij}, \beta_{ij})$，其逆为$A_{(E_j, E_i)} = (-\alpha_{ij}/\beta_{ij}, 1/\beta_{ij})$（当$\beta_{ij} \neq 0$时）。

2-形式场强（审计通量）$F = \dd_1 A$在2-单形$\Sigma = (E_i, E_j, E_k)$上定义为校准模型组合的\textbf{非传递性}（non-transitivity）：

\begin{equation}\label{eq:F_def}
    F_{ijk} = \bigl\| A_{(E_i, E_j)} \circ A_{(E_j, E_k)} - A_{(E_i, E_k)}^{-1} \bigr\|,
\end{equation}

其中$\circ$表示校准模型的函数复合：若$S_i = \alpha_{ij} + \beta_{ij} S_j$且$S_j = \alpha_{jk} + \beta_{jk} S_k$，则复合模型为$S_i = (\alpha_{ij} + \beta_{ij}\alpha_{jk}) + (\beta_{ij}\beta_{jk}) S_k$。

\textbf{物理意义：} $F_{ijk}$度量了通过专家$E_j$间接校准$E_i \to E_k$与直接校准$E_i \to E_k$之间的不一致程度。若所有两两校准模型可传递（即$F_{ijk}=0$对所有的$(i,j,k)$），则存在全局一致的参考框架。

若专家共享同一底层评分框架（仅常数偏移），则$\alpha_{ij} = c_i - c_j$和$\beta_{ij} = 1$对所有$(i,j)$成立。此时$F_{ijk} = 0$——校准可传递，对应平坦联络。
\end{definition}

\begin{remark}[1-循环瞬子被杀死，但校准非传递性存活 / 1-Cycle Killed, Calibration Non-Transitivity Survives]
\label{rem:1cycle_death}
如果每个专家的评分都是同一底层真值加上常数偏移：$S_i(x) = S^*(x) + c_i$，则两两校准模型具有简单形式：$\alpha_{ij} = c_i - c_j$，$\beta_{ij} = 1$。此时校准模型的复合是传递的：
\[
(\alpha_{ij} + \beta_{ij}\alpha_{jk}, \beta_{ij}\beta_{jk}) = (c_i - c_j + 1 \cdot (c_j - c_k), 1 \cdot 1) = (c_i - c_k, 1),
\]
而直接校准$S_k \to S_i$（逆模型）也为$(\alpha_{ik}, \beta_{ik}) = (c_i - c_k, 1)$。两者一致，因此$F_{ijk} = 0$。

这就是C7结论的推广：\textbf{1-循环瞬子被杀死}——当所有专家共享同一线性评分框架时，不存在非平凡审计通量。$A$是恰当形式：存在全局0-上链$\phi$（即全局参考评分$S^* + c_i$）使$A$可由此导出。

然而，当专家评分来自\textbf{不同评估框架}时——即$\beta_{ij} \neq 1$，或评分之间关系非线性——校准模型的复合不再传递。这正是C8的新物理：\textbf{2-形式通量在专家评分框架非线性/非同构时存活}。
\end{remark}

\bigskip
\noindent
\textbf{Remark (The 1-Cycle Instanton Cannot Survive).}
The above shows: if $F$ is computed directly via substitution from \eqref{eq:F_expanded}, $F_\Sigma$ is identically zero. This means any 1-form $A$ defined by \eqref{eq:A_def} satisfies $\dd_1 A = 0$ (is a cocycle). C7 inferred from this that the 1-cycle audit instanton is killed — because $A$ is an exact form, i.e., there exists a 0-cochain $S$ such that $A = \dd_0 S$.

However, this rests on the default assumption that ``all experts share the same Cercis scoring function.'' In a real audit, experts may have \textbf{substantive disagreements} on the same data point $x$ — their scores $S_i(x)$ are not merely noisy perturbations of the same ground-truth function, but reflect fundamentally different evaluative frameworks. In this case, the connection $A$ is no longer derived from a single global $S^*$, and $\dd_1 A$ can be non-zero.

The next section proves: the existence of disagreeing experts is equivalent to the co-exact part of $A$ being non-zero, which in turn is equivalent to non-zero 2-form flux.

% ============================================================================
\section{1-循环审计瞬子的死亡 / Death of the 1-Cycle Audit Instanton}
\label{sec:death}
% ============================================================================

\noindent
\textbf{中文.}
在详细展开高维瞬子理论之前，我们首先严格回顾C7的结论，作为必要的理论背景。

\begin{theorem}[C7: 1-循环瞬子死亡定理 / 1-Cycle Instanton Death Theorem]
\label{thm:C7}
设$S_i(x) = S^*(x) + c_i$对所有专家$E_i$和数据点$x$成立（即所有专家共享同一底层评分函数，仅常数偏移不同）。则两两校准模型为$A_{(E_i, E_j)} = (c_i - c_j, 1)$，且：

\begin{enumerate}[label=(\roman*), leftmargin=*]
    \item $F_{ijk} = 0$对所有专家三元组成立（场强为零，平坦联络）；
    \item 校准复合是传递的（transitive）：$A_{(E_i, E_j)} \circ A_{(E_j, E_k)} = A_{(E_i, E_k)}^{-1}$；
    \item 存在全局参考框架：所有专家的评分可通过常数偏移$c_i$规范到同一基底。
\end{enumerate}

其逆否命题同样重要：若检测到$F_{ijk} \neq 0$，则专家评分不能表示为$S_i(x) = S^*(x) + c_i$——即不存在全局一致的线性参考框架。
\end{theorem}

\begin{proof}
(i) $\dd_1 A = \dd_1(\dd_0 \phi) = (\dd_1 \circ \dd_0)\phi = 0$，因为上边缘算子的平方为零。

(ii) 对于任一闭路$\gamma$（可分解为1-单形链），$\hol_\gamma(A) = \sum_{\sigma \in \gamma} A_\sigma$。由于$A = \dd_0 \phi$且在闭路上$\dd_0 \phi$的求和项相互抵消（类似于微积分基本定理），和乐为零。

(iii) 规范变换$\delta A = \dd_0 \chi$导致$\delta F = \dd_1(\delta A) = \dd_1(\dd_0 \chi) = 0$。$\hfill\square$
\end{proof}

\begin{proof}[English Proof Sketch]
(i) $\dd_1 A = \dd_1(\dd_0 \phi) = (\dd_1 \circ \dd_0)\phi = 0$, since the coboundary operator squares to zero.

(ii) For any closed path $\gamma$ (decomposable into a chain of 1-simplices), $\hol_\gamma(A) = \sum_{\sigma \in \gamma} A_\sigma$. Since $A = \dd_0 \phi$, terms of $\dd_0 \phi$ cancel pairwise along the closed loop (analogous to the fundamental theorem of calculus), yielding zero holonomy.

(iii) The gauge transformation $\delta A = \dd_0 \chi$ induces $\delta F = \dd_1(\delta A) = \dd_1(\dd_0 \chi) = 0$. $\hfill\square$
\end{proof}

\begin{corollary}[全局一致性检测 / Global Consistency Detection]
\label{cor:C7_corollary}
若对于所有2-单形$\Sigma$均有$F_\Sigma = 0$，则存在全局0-上链$\phi$（即``真实Cercis评分函数''）使得所有专家的评分可表示为$S_i(x) = \phi(x) + c_i$，其中$c_i$是仅依赖于专家的常数偏移。这意味着所有不一致性均为\textbf{可规范的}（gauge-able）。
\end{corollary}

% ============================================================================
\section{高维审计瞬子：$k>1$的非平坦构型 / Higher-Dimensional Audit Instantons: Non-Flat Configurations at $k>1$}
\label{sec:nonflat}
% ============================================================================

\subsection{超越线性框架：非平坦联络的产生}
\label{sec:beyond_exact}

\noindent
\textbf{中文.}
C7的``瞬子死亡''结论成立的前提是：所有专家评分可表示为$S_i(x) = S^*(x) + c_i$（同一底层函数加常数偏移）。在此前提下，两两校准模型为$\alpha_{ij} = c_i - c_j$，$\beta_{ij} = 1$，复合可传递。

但在真实审计中，专家可能使用\textbf{根本不同的评估框架}：
\begin{itemize}
    \item \textbf{非线性尺度：} 专家$E_i$可能使用压缩尺度（如$\sqrt{S^*}$），而$E_j$使用扩展尺度（如$(S^*)^2$）。此时$\beta_{ij} \neq 1$，校准模型复合不再传递——产生非零$F$。
    \item \textbf{异质评估维度：} 专家可能强调数据的不同方面（如语法 vs.\ 语义、局部 vs.\ 全局特征），导致评分之间的关系不能简约到单一线性变换。
    \item \textbf{非单调响应：} 某些专家可能对高质量和低质量样本有不同的敏感度，导致评分映射不具有全局线性性。
\end{itemize}

在这些情况下，即使每对专家之间的线性回归可以很好地拟合（高$R^2$），三个模型的复合仍可能与直接模型冲突——这就是\textbf{校准非传递性}，即非平坦2-形式通量。

\vspace{0.5em}
\noindent
\textbf{English.}
C7's ``instanton death'' conclusion holds under the premise that all expert scores can be expressed as $S_i(x) = S^*(x) + c_i$ (same underlying function plus constant offset). Under this premise, pairwise calibration models are $\alpha_{ij} = c_i - c_j$, $\beta_{ij} = 1$, and composition is transitive.

But in real audits, experts may use \textbf{fundamentally different evaluative frameworks}:
\begin{itemize}
    \item \textbf{Nonlinear scales:} Expert $E_i$ may use a compressive scale (e.g., $\sqrt{S^*}$) while $E_j$ uses an expansive scale (e.g., $(S^*)^2$). Then $\beta_{ij} \neq 1$, and calibration model composition is no longer transitive — producing non-zero $F$.
    \item \textbf{Heterogeneous evaluation dimensions:} Experts may emphasize different aspects of data (e.g., syntax vs.\ semantics, local vs.\ global features), such that score relationships cannot be reduced to a single linear transformation.
    \item \textbf{Non-monotonic responses:} Some experts may have different sensitivity to high-quality vs.\ low-quality samples, making score mappings not globally linear.
\end{itemize}

In these cases, even though each pairwise linear regression may fit well (high $R^2$), the composition of three models may conflict with the direct model — this is \textbf{calibration non-transitivity}, i.e., non-flat 2-form flux.

% -------------------------------------------------------------------
\subsection{Hodge分解定理}
\label{sec:hodge}

\begin{definition}[组合Laplace算子 / Combinatorial Laplacian]
\label{def:laplacian}
在有限单纯复形$\SSS$上，$k$-阶组合Laplace算子$\Delta_k: C^k \to C^k$定义为：

\begin{equation}\label{eq:laplacian}
    \Delta_k = \dd_{k-1} \circ \delta_{k-1} + \delta_k \circ \dd_k,
\end{equation}

其中$\delta_k: C^{k+1} \to C^k$是上边缘算子$\dd_k$关于标准内积的伴随算子（adjoint）。具体地，在标准正交基下，$\delta_k$由$\dd_k$的转置矩阵给出。对于$k=0$，第一项消失（$\dd_{-1} = 0$）；对于$k = \dim \SSS$，第二项消失。
\end{definition}

\bigskip
\noindent
\textbf{Definition (Combinatorial Laplacian).}
On a finite simplicial complex $\SSS$, the $k$-th combinatorial Laplacian $\Delta_k: C^k \to C^k$ is defined as:

\begin{equation}
    \Delta_k = \dd_{k-1} \circ \delta_{k-1} + \delta_k \circ \dd_k,
\end{equation}

where $\delta_k: C^{k+1} \to C^k$ is the adjoint of $\dd_k$ with respect to the standard inner product. Concretely, in the standard orthonormal basis, $\delta_k$ is given by the transpose of $\dd_k$. For $k=0$, the first term vanishes ($\dd_{-1} = 0$); for $k = \dim \SSS$, the second term vanishes.

\begin{theorem}[Hodge分解定理 / Hodge Decomposition Theorem]
\label{thm:hodge}
对于有限单纯复形$\SSS$上的$k$-上链空间$C^k(\SSS; \R)$，存在正交直和分解：

\begin{equation}\label{eq:hodge_dec}
    C^k(\SSS; \R) = \im \dd_{k-1} \;\oplus\; \im \delta_k \;\oplus\; \ker \Delta_k.
\end{equation}

这三个子空间分别称为：
\begin{itemize}
    \item $\im \dd_{k-1}$: \textbf{恰当$k$-上链} (exact $k$-cochains) — 可表示为低一维上链的上边缘；
    \item $\im \delta_k$: \textbf{共存$k$-上链} (co-exact $k$-cochains) — 可表示为高一维上链的余边缘（codifferential）；
    \item $\ker \Delta_k$: \textbf{调和$k$-上链} (harmonic $k$-cochains) — 满足$\Delta_k \omega = 0$，等价于同时满足$\dd_k\omega = 0$和$\delta_{k-1}\omega = 0$。与上同调群$H^k(\SSS; \R)$同构。
\end{itemize}
\end{theorem}

\begin{proof}
这是有限维向量空间上椭圆算子理论的标准结果。关键步骤：

\textbf{步骤1:} 证明$\im \dd_{k-1} \perp \im \delta_k$。对于任意$\alpha \in C^{k-1}$和$\beta \in C^{k+1}$：
\[
\inner{\dd_{k-1}\alpha}{\delta_k\beta} = \inner{\alpha}{\delta_{k-1}\delta_k\beta} = \inner{\alpha}{0} = 0,
\]
因为$\delta_{k-1}\delta_k = (\dd_k \dd_{k-1})^* = 0^* = 0$。

\textbf{步骤2:} 证明$\im \dd_{k-1} \perp \ker \Delta_k$。对于任意$\alpha \in C^{k-1}$和$\omega \in \ker \Delta_k$：
\[
\inner{\dd_{k-1}\alpha}{\omega} = \inner{\alpha}{\delta_{k-1}\omega}.
\]
由于$\omega \in \ker \Delta_k$，有$\delta_{k-1}\omega = 0$（因为$\Delta_k\omega = \dd_{k-1}\delta_{k-1}\omega + \delta_k\dd_k\omega = 0$，且两项均非负，故每一项必须为零）。因此内积为零。

\textbf{步骤3:} 类似可证$\im \delta_k \perp \ker \Delta_k$。

\textbf{步骤4:} 直和性质：$\im \dd_{k-1} \oplus \im \delta_k = (\ker \Delta_k)^\perp$。因为$\Delta_k$是自伴算子，其核的正交补等于其像。$\hfill\square$
\end{proof}

\begin{proof}[English Proof Sketch]
Standard result from elliptic operator theory on finite-dimensional vector spaces.

\textbf{Step 1:} Show $\im \dd_{k-1} \perp \im \delta_k$. For any $\alpha \in C^{k-1}$, $\beta \in C^{k+1}$:
\[
\inner{\dd_{k-1}\alpha}{\delta_k\beta} = \inner{\alpha}{\delta_{k-1}\delta_k\beta} = \inner{\alpha}{0} = 0,
\]
since $\delta_{k-1}\delta_k = (\dd_k \dd_{k-1})^* = 0^* = 0$.

\textbf{Step 2:} Show $\im \dd_{k-1} \perp \ker \Delta_k$. For any $\alpha \in C^{k-1}$, $\omega \in \ker \Delta_k$:
\[
\inner{\dd_{k-1}\alpha}{\omega} = \inner{\alpha}{\delta_{k-1}\omega}.
\]
Since $\omega \in \ker \Delta_k$, we have $\delta_{k-1}\omega = 0$ (because $\Delta_k\omega = \dd_{k-1}\delta_{k-1}\omega + \delta_k\dd_k\omega = 0$ and both terms are non-negative, so each must vanish). Hence the inner product is zero.

\textbf{Step 3:} $\im \delta_k \perp \ker \Delta_k$ is analogous.

\textbf{Step 4:} The direct sum property: $\im \dd_{k-1} \oplus \im \delta_k = (\ker \Delta_k)^\perp$, because $\Delta_k$ is self-adjoint and the orthogonal complement of its kernel equals its image. $\hfill\square$
\end{proof}

% -------------------------------------------------------------------
\subsection{场强与共存部分的关系}
\label{sec:F_coexact}

\begin{theorem}[场强-共存关系定理 / Field-Strength Co-Exact Relation]
\label{thm:F_coexact}
设$A \in C^1(\SSS; \R)$为1-形式审计联络，其Hodge分解为：

\begin{equation}\label{eq:A_hodge}
    A = A_{\mathrm{exact}} + A_{\mathrm{coexact}} + A_{\mathrm{harmonic}},
\end{equation}

其中$A_{\mathrm{exact}} \in \im \dd_0$，$A_{\mathrm{coexact}} \in \im \delta_1$，$A_{\mathrm{harmonic}} \in \ker \Delta_1$。

则2-形式场强$F = \dd_1 A$满足：

\begin{equation}\label{eq:F_from_coexact}
    F = \dd_1 A_{\mathrm{coexact}}.
\end{equation}

特别地：
\begin{enumerate}[label=(\roman*), leftmargin=*]
    \item $F = 0$ 当且仅当 $A_{\mathrm{coexact}} = 0$（在$\dd_1$为单射的子空间上）；
    \item $A_{\mathrm{exact}}$和$A_{\mathrm{harmonic}}$对$F$无贡献；
    \item $F_{\mathrm{norm}} \equiv \norm{F}_2 = \norm{\dd_1 A_{\mathrm{coexact}}}_2$，可用作共存部分大小的度量。
\end{enumerate}
\end{theorem}

\begin{proof}
将Hodge分解代入$F = \dd_1 A$：

\[
F = \dd_1(A_{\mathrm{exact}} + A_{\mathrm{coexact}} + A_{\mathrm{harmonic}}).
\]

逐项分析：
\begin{itemize}
    \item $\dd_1 A_{\mathrm{exact}} = \dd_1(\dd_0 \phi) = 0$（对于某$\phi \in C^0$），因为$\dd_1 \circ \dd_0 = 0$。
    \item $\dd_1 A_{\mathrm{harmonic}} = 0$，因为$A_{\mathrm{harmonic}} \in \ker \Delta_1$意味着$\dd_1 A_{\mathrm{harmonic}} = 0$（调和形式是上闭链）。
    \item 因此$F = \dd_1 A_{\mathrm{coexact}}$。
\end{itemize}

(i) 若$A_{\mathrm{coexact}} = 0$，显然$F = 0$。反之，若$F = 0$，则$\dd_1 A_{\mathrm{coexact}} = 0$，即$A_{\mathrm{coexact}} \in \ker \dd_1$。但$A_{\mathrm{coexact}} \in \im \delta_1$，而$\im \delta_1 \cap \ker \dd_1 = \{0\}$（因为$\im \delta_1 \subseteq (\ker \dd_1)^\perp$，来自Hodge分解的正交性）。因此$A_{\mathrm{coexact}} = 0$。

(ii) 由上直接得出。

(iii) $F_{\mathrm{norm}}$是$\dd_1$限制在$\im \delta_1$上的算子范数，由于$\dd_1$在有限维空间上是有界算子，$\norm{F}_2$正比于$\norm{A_{\mathrm{coexact}}}_2$（在紧致复形上，通过离散Poincaré不等式）。$\hfill\square$
\end{proof}

\begin{proof}[English Proof Sketch]
Substitute the Hodge decomposition into $F = \dd_1 A$:

\[
F = \dd_1(A_{\mathrm{exact}} + A_{\mathrm{coexact}} + A_{\mathrm{harmonic}}).
\]

Term-by-term:
\begin{itemize}
    \item $\dd_1 A_{\mathrm{exact}} = \dd_1(\dd_0 \phi) = 0$ (for some $\phi \in C^0$), because $\dd_1 \circ \dd_0 = 0$.
    \item $\dd_1 A_{\mathrm{harmonic}} = 0$, because $A_{\mathrm{harmonic}} \in \ker \Delta_1$ implies $\dd_1 A_{\mathrm{harmonic}} = 0$ (harmonic forms are cocycles).
    \item Hence $F = \dd_1 A_{\mathrm{coexact}}$.
\end{itemize}

(i) If $A_{\mathrm{coexact}} = 0$, clearly $F = 0$. Conversely, if $F = 0$, then $\dd_1 A_{\mathrm{coexact}} = 0$, i.e., $A_{\mathrm{coexact}} \in \ker \dd_1$. But $A_{\mathrm{coexact}} \in \im \delta_1$, and $\im \delta_1 \cap \ker \dd_1 = \{0\}$ (from orthogonality of the Hodge decomposition). Hence $A_{\mathrm{coexact}} = 0$.

(ii) Follows directly from the above.

(iii) $F_{\mathrm{norm}}$ is the operator norm of $\dd_1$ restricted to $\im \delta_1$; since $\dd_1$ is a bounded operator on a finite-dimensional space, $\norm{F}_2$ is proportional to $\norm{A_{\mathrm{coexact}}}_2$ (via the discrete Poincaré inequality on compact complexes). $\hfill\square$
\end{proof}

\begin{corollary}[审计不一致性的谱特征 / Spectral Signature of Audit Inconsistency]
\label{cor:spectral}
将$F_{\mathrm{norm}}$视为审计不一致性的全局度量。对于$N$个专家在$M$个数据点上的审计，定义不一致性矩阵$\mathbf{D} \in \R^{N \times N}$：

\begin{equation}\label{eq:D_matrix}
    \mathbf{D}_{ij} = \frac{1}{M} \sum_{x \in \mathcal{X}} |S_i(x) - S_j(x)|.
\end{equation}

则以下条件是等价的：
\begin{enumerate}[label=(\roman*), leftmargin=*]
    \item 所有不一致性可规范（gauge-able），即存在偏移$\{c_i\}$使得$S_i(x) - c_i = S_j(x) - c_j$对所有$i,j,x$成立；
    \item $F_\Sigma = 0$对所有2-单形$\Sigma$成立；
    \item $\mathbf{D}$满足三角不等式$\mathbf{D}_{ik} \leq \mathbf{D}_{ij} + \mathbf{D}_{jk}$且$\mathbf{D}_{ii} = 0$，即$\mathbf{D}$是一个度量；
    \item Hodge分解中$A_{\mathrm{coexact}} = 0$。
\end{enumerate}
\end{corollary}

% ============================================================================
\section{Situs复形上的2-闭链与审计通量 / 2-Cycles on the Situs Complex and Audit Flux}
\label{sec:2cycles}
% ============================================================================

\subsection{专家图上的2-闭链结构}
\label{sec:2cycle_construction}

\noindent
\textbf{中文.}
在更新后的Situs复形中，顶点为专家$\{E_i\}$，边为专家对$(E_i, E_j)$，校准联络$A$在边上取值$(\alpha_{ij}, \beta_{ij})$。2-单形为专家三元组$(E_i, E_j, E_k)$，场强$F_{ijk}$衡量校准非传递性。

2-闭链的构造在此框架下变得自然：

\begin{definition}[专家图上的2-闭链与通量 / 2-Cycle and Flux on Expert Graph]
\label{def:2cycle_flux}
设$K_N$为$N$位专家的完全图。2-链$\Sigma$是三角形的整系数形式和：
\begin{equation}\label{eq:2chain}
    \Sigma = \sum_{1 \leq i < j < k \leq N} c_{ijk} \cdot \Delta_{ijk},
\end{equation}
其中$\Delta_{ijk}$为三角形$(E_i, E_j, E_k)$，$c_{ijk} \in \Z$。$\Sigma$是2-闭链当且仅当$\bd_2 \Sigma = 0$。

审计通量积分定义为：
\begin{equation}\label{eq:flux_integral}
    \Phi(\Sigma) = \sum_{i<j<k} c_{ijk} \cdot F_{ijk}.
\end{equation}

当完全图$K_N$（$N \geq 3$）中检测到任意三角形上的$F_{ijk} \neq 0$时，即存在非平凡审计通量。
\end{definition}

\subsection{校准非传递性的代数判据}
\label{sec:algebraic_criterion}

\begin{theorem}[校准非传递性判据 / Calibration Non-Transitivity Criterion]
\label{thm:2cycle_computation}
对于三位专家$E_i, E_j, E_k$，校准非传递性$F_{ijk}$的计算涉及以下分量：

\begin{equation}\label{eq:F_components}
    F_{ijk}^\alpha = |\alpha_{ij} + \beta_{ij}\alpha_{jk} - (-\alpha_{ki}/\beta_{ki})|, \quad
    F_{ijk}^\beta = |\beta_{ij}\beta_{jk} - 1/\beta_{ki}|.
\end{equation}

总场强大小$F_{ijk} = \sqrt{(F_{ijk}^\alpha)^2 + (F_{ijk}^\beta)^2}$。

\textbf{几何解释：} $F_{ijk}^\beta$度量尺度的非传递性（$\beta$之积偏离1），$F_{ijk}^\alpha$度量偏移的非传递性（校准后的截距不一致）。当所有$\beta = 1$且$\alpha_{ij} + \alpha_{jk} = \alpha_{ik}$时，恢复平坦联络。
\end{theorem}

\begin{proof}
直接计算：复合校准$(E_i \to E_j) \circ (E_j \to E_k)$给出$S_i = (\alpha_{ij} + \beta_{ij}\alpha_{jk}) + (\beta_{ij}\beta_{jk}) S_k$。逆校准$(E_k \to E_i)$的逆给出$S_i = -\alpha_{ki}/\beta_{ki} + (1/\beta_{ki}) S_k$。差异范数即为$F_{ijk}$。$\hfill\square$
\end{proof}

\subsection{非传递性的信息论度量}
\label{sec:info_measure}

\begin{definition}[循环不一致性信噪比 / Circular Inconsistency SNR]
\label{def:circ_snr}
对于三角形$(E_i, E_j, E_k)$，定义循环不一致性的信噪比：

\begin{equation}\label{eq:circ_snr_def}
    \mathrm{SNR}_{ijk} = \frac{F_{ijk}}{\sigma_{\mathrm{noise}}},
\end{equation}

其中$\sigma_{\mathrm{noise}}$为两两校准残差的平均标准差。高SNR表示检测到的非传递性远大于估计噪声，构成统计显著的循环不一致性。
\end{definition}

% ============================================================================
\section{物理意义：循环不一致性 / Physical Meaning: Circular Inconsistency}
\label{sec:physical}
% ============================================================================

\subsection{三角不一致性的信息论模型}
\label{sec:triangular}

\noindent
\textbf{中文.}
审计通量$F_{ijk} \neq 0$的物理含义：三位专家的两两校准模型不能通过任意单一专家的参考框架统一。具体地，如果通过$E_j$间接校准$E_i \to E_k$的结果与直接校准$E_i \to E_k$的结果不一致，则三位专家的评分体系构成\textbf{循环不一致性}。

考虑三位专家$A, B, C$评估同一组数据。可能发生以下情况：
\begin{itemize}
    \item $A$和$B$之间的线性校准很好（高$R^2$）；
    \item $B$和$C$之间的线性校准也很好；
    \item 但通过$B$间接校准$A \to C$的结果与直接校准$A \to C$显著不同（$F_{ABC} \gg 0$）。
\end{itemize}

此时，尽管每对专家之间的线性模型拟合良好，三者的校准体系全局上却是自相矛盾的。这就是\textbf{校准非传递性}——循环不一致性的操作化定义。

\textbf{实际审计示例：} 三位审稿人评估论文。审稿人A使用严格标准（分数压缩在3-5分），审稿人B使用宽松标准（1-5分全范围），审稿人C对理论文章宽松但对实验文章严格。任何一对之间的线性校准都能拟合（因为有共同的文章子集），但三条校准曲线不能闭合为一致的全局框架。

\vspace{0.5em}
\noindent
\textbf{English.}
The physical meaning of audit flux $\Phi(\Sigma) \neq 0$ is profound: it signals \textbf{circular inconsistency} — a systematic inconsistency among three or more experts that cannot be reduced to simple pairwise disagreement.

Consider three experts $A, B, C$ evaluating the same set of data. The following can occur:
\begin{itemize}
    \item $A$ and $B$ agree in their evaluative framework (their score discrepancies are eliminable via constant offset);
    \item $B$ and $C$ also agree in their evaluative framework;
    \item But $A$ and $C$'s evaluative frameworks conflict — cannot be reconciled via constant offset.
\end{itemize}

In this case, although each pair of experts is ``apparently consistent'' (1-form holonomy zero), globally the three are self-contradictory. This is \textbf{circular inconsistency}.

% -------------------------------------------------------------------
\subsection{形式化模型}
\label{sec:formal_model}

\begin{definition}[循环不一致性 / Circular Inconsistency]
\label{def:circular}
设$\mathcal{E} = \{E_1, \ldots, E_N\}$为$N \geq 3$位专家。专家$E_i$和$E_j$之间称为\textbf{两两一致}（pairwise consistent），如果存在常数$c_{ij} \in \R$使得对所有数据点$x \in \mathcal{X}$：

\begin{equation}\label{eq:pairwise_cons}
    |S_i(x) - S_j(x) - c_{ij}| \leq \varepsilon,
\end{equation}

其中$\varepsilon \geq 0$为容许公差。

专家集合$\mathcal{E}$称为\textbf{全局一致}（globally consistent），如果存在一组常数$\{c_i\}_{i=1}^N$使得对所有$i,j$：

\begin{equation}\label{eq:global_cons}
    c_{ij} = c_i - c_j.
\end{equation}

如果所有专家对都是两两一致的，但不是全局一致的，则称存在\textbf{循环不一致性}。
\end{definition}

\bigskip
\noindent
\textbf{Definition (Circular Inconsistency).}
Let $\mathcal{E} = \{E_1, \ldots, E_N\}$ be $N \geq 3$ experts. Experts $E_i$ and $E_j$ are called \textbf{pairwise consistent} if there exists a constant $c_{ij} \in \R$ such that for all data points $x \in \mathcal{X}$:

\begin{equation}
    |S_i(x) - S_j(x) - c_{ij}| \leq \varepsilon,
\end{equation}

where $\varepsilon \geq 0$ is the tolerance.

The expert set $\mathcal{E}$ is called \textbf{globally consistent} if there exist constants $\{c_i\}_{i=1}^N$ such that for all $i,j$:

\begin{equation}
    c_{ij} = c_i - c_j.
\end{equation}

If all expert pairs are pairwise consistent but not globally consistent, we say there is \textbf{circular inconsistency}.

\begin{theorem}[循环不一致性 $\iff$ 非零审计通量 / Circular Inconsistency $\iff$ Non-Zero Audit Flux]
\label{thm:circular_flux}
对于$N \geq 3$位专家，以下条件等价：

\begin{enumerate}[label=(\roman*), leftmargin=*]
    \item 存在循环不一致性；
    \item 存在2-闭链$\Sigma$使得审计通量$\Phi(\Sigma) \neq 0$；
    \item 偏移矩阵$\mathbf{C} = (c_{ij})$不满足上边缘条件：不存在向量$\mathbf{c} = (c_1, \ldots, c_N)$使得$c_{ij} = c_j - c_i$对所有$i,j$成立；
    \item $\dd_1 \mathbf{C} \neq 0$，其中$\dd_1$作用于以专家为顶点的完全图上的1-上链。
\end{enumerate}
\end{theorem}

\begin{proof}
(i) $\iff$ (iii)：直接来自定义。

(iii) $\iff$ (iv)：在完全图$K_N$上，$\mathbf{C}$可视为1-上链（在边$(i,j)$上取值$c_{ij}$）。$\dd_1 \mathbf{C} = 0$当且仅当$\mathbf{C}$是上边缘，即存在0-上链$\mathbf{c}$使得$\mathbf{C} = \dd_0 \mathbf{c}$。这正是条件$c_{ij} = c_j - c_i$。

(ii) $\iff$ (iv)：审计通量$F = \dd_1 A$，其中$A$在边$(i,j,x)$上取值为$S_i(x) - S_j(x)$。由于$\Phi(\Sigma) = \int_\Sigma F = \int_\Sigma \dd_1 A$，且Stokes定理给出$\Phi(\Sigma) = \oint_{\bd\Sigma} A = 0$（因为$\Sigma$是闭链，$\bd\Sigma = 0$）——但这是指\textbf{几何}积分。在\textbf{组合}语境下，$\Phi(\Sigma) \neq 0$当且仅当$\dd_1 A$作为2-上链非零，这要求在某个三角形$(i,j,k)$上$A_{ij} + A_{jk} + A_{ki} \neq 0$。

现在，如果我们将$A$与偏移矩阵$\mathbf{C}$的适当嵌入相关联，可以看到：$A_{ij} + A_{jk} + A_{ki} \neq 0$当且仅当$c_{ij} + c_{jk} + c_{ki} \neq 0$（在数据平均意义上），这恰好是$\dd_1 \mathbf{C} \neq 0$的条件。$\hfill\square$
\end{proof}

\begin{proof}[English Proof Sketch]
(i) $\iff$ (iii): Direct from definition.

(iii) $\iff$ (iv): On the complete graph $K_N$, $\mathbf{C}$ can be viewed as a 1-cochain (taking value $c_{ij}$ on edge $(i,j)$). $\dd_1 \mathbf{C} = 0$ iff $\mathbf{C}$ is a coboundary, i.e., there exists a 0-cochain $\mathbf{c}$ such that $\mathbf{C} = \dd_0 \mathbf{c}$. This is precisely the condition $c_{ij} = c_j - c_i$.

(ii) $\iff$ (iv): The audit flux $F = \dd_1 A$, where $A$ on edge $(i,j,x)$ takes value $S_i(x) - S_j(x)$. $\Phi(\Sigma) \neq 0$ iff $\dd_1 A$ is non-zero as a 2-cochain, which requires that on some triangle $(i,j,k)$, $A_{ij} + A_{jk} + A_{ki} \neq 0$.

Now, relating $A$ to an appropriate embedding of the offset matrix $\mathbf{C}$, we see that $A_{ij} + A_{jk} + A_{ki} \neq 0$ iff $c_{ij} + c_{jk} + c_{ki} \neq 0$ (in the data-averaged sense), which is precisely the condition $\dd_1 \mathbf{C} \neq 0$. $\hfill\square$
\end{proof}

% -------------------------------------------------------------------
\subsection{信息论下界}
\label{sec:info_bound}

\begin{theorem}[循环不一致性的信息论下界 / Information-Theoretic Lower Bound for Circular Inconsistency]
\label{thm:info_bound}
设三位专家$A, B, C$的两两Cercis评分差异具有方差$\sigma_{AB}^2, \sigma_{BC}^2, \sigma_{CA}^2$。定义循环不一致性的信噪比：

\begin{equation}\label{eq:SNR}
    \mathrm{SNR}_{\mathrm{circ}} = \frac{|\E[F_{ABC}]|^2}{\Var[F_{ABC}]},
\end{equation}

其中$F_{ABC} = (S_A - S_B) + (S_B - S_C) + (S_C - S_A)$在一个数据点上的取值。

若评分差异在两两之间相互独立，则：

\begin{equation}\label{eq:SNR_bound}
    \Var[F_{ABC}] = \sigma_{AB}^2 + \sigma_{BC}^2 + \sigma_{CA}^2.
\end{equation}

要在显著水平$\alpha$下检测循环不一致性（即拒绝$F_{ABC} = 0$的零假设），所需数据点数量$M$满足：

\begin{equation}\label{eq:M_bound}
    M \geq \frac{(\sigma_{AB}^2 + \sigma_{BC}^2 + \sigma_{CA}^2) \cdot (z_{1-\alpha/2})^2}{\delta^2},
\end{equation}

其中$\delta = |\E[F_{ABC}]|$为期望循环不一致性强度，$z_{1-\alpha/2}$为标准正态分位数。
\end{theorem}

\begin{proof}
在独立性假设下，$F_{ABC}$的方差为三对独立差异的方差之和。由于$F_{ABC}$是$M$个独立同分布观测的均值，其样本均值的方差为$\Var[F_{ABC}]/M$。在中心极限定理下，使用标准正态近似得到Wald检验的样本量公式。$\hfill\square$
\end{proof}

\begin{proof}[English Proof]
Under independence, the variance of $F_{ABC}$ is the sum of variances of the three independent pairwise differences. Since $F_{ABC}$ is the mean of $M$ i.i.d.\ observations, the variance of the sample mean is $\Var[F_{ABC}]/M$. Under the central limit theorem, the standard normal approximation yields the Wald-test sample size formula. $\hfill\square$
\end{proof}

% ============================================================================
\section{检测算法 / Detection Algorithm}
\label{sec:algorithm}
% ============================================================================

\noindent
\textbf{中文.}
本节描述在Situs复形的2-骨架上检测非零审计通量的算法。算法核心思想：扫描所有可能的2-闭链，计算其通量积分，标记超出阈值$\tau$的2-闭链。

\vspace{0.5em}
\noindent
\textbf{English.}
This section describes the algorithm for detecting non-zero audit flux on the 2-skeleton of the Situs complex. Core idea: scan all possible 2-cycles, compute their flux integrals, and flag those exceeding threshold $\tau$.

% -------------------------------------------------------------------
\subsection{算法描述}
\label{sec:algo_desc}

\begin{algorithm}[h]
\caption{审计通量检测算法 / Audit Flux Detection Algorithm}
\label{alg:flux_detect}
\begin{algorithmic}[1]
\Require 专家评分矩阵 $\mathbf{S} \in \R^{N \times M}$ ($N$位专家，$M$个数据点)
\Require 检测阈值 $\tau > 0$
\Ensure 不一致2-闭链的标记集合 $\mathcal{F}_{\mathrm{bad}}$

\State \textbf{步骤 1:} 构造Situs复形的2-骨架
\For{$x = 1$ \textbf{to} $M$}
    \For{each triple $(i,j,k)$ with $1 \leq i < j < k \leq N$}
        \State 添加2-单形 $\sigma_{x,ijk}$ 到2-骨架 $K_2$
    \EndFor
\EndFor

\State \textbf{步骤 2:} 计算2-形式场强
\For{each 2-单形 $\sigma = (x, i, j, k) \in K_2$}
    \State $F_{x,ijk} \gets S_i(x) - S_j(x) + S_j(x) - S_k(x) + S_k(x) - S_i(x)$
    \Comment{注意：在单点三角上恒为零；非零仅出现在跨点2-闭链}
\EndFor

\State \textbf{步骤 3:} 枚举跨点2-闭链
\State 初始化 $\mathcal{F}_{\mathrm{bad}} \gets \emptyset$
\For{each expert triple $(i,j,k)$}
    \For{each triple of data points $(x_1, x_2, x_3)$ with $x_1 < x_2 < x_3$}
        \State 构造跨点2-闭链 $\Sigma = \sigma_{x_1,ijk} - \sigma_{x_2,ijk} + \sigma_{x_3,ijk}$
        \State 验证 $\bd_2 \Sigma = 0$（检查边界）
        \If{$\bd_2 \Sigma = 0$}
            \State $\Phi \gets F_{x_1,ijk} - F_{x_2,ijk} + F_{x_3,ijk}$
            \If{$|\Phi| > \tau$}
                \State $\mathcal{F}_{\mathrm{bad}} \gets \mathcal{F}_{\mathrm{bad}} \cup \{(\Sigma, \Phi)\}$
            \EndIf
        \EndIf
    \EndFor
\EndFor

\State \textbf{步骤 4:} 计算跨点通量的替代度量（更高效）
\State 构造差异矩阵 $\mathbf{D} \in \R^{M \times N \times N}$：
\For{$x = 1$ \textbf{to} $M$}
    \For{$i,j = 1$ \textbf{to} $N$}
        \State $\mathbf{D}_{x,ij} \gets |S_i(x) - S_j(x)|$
    \EndFor
\EndFor

\State 对每个专家三元组$(i,j,k)$，计算：
\State $\mathrm{FluxScore}_{ijk} \gets \frac{1}{M}\sum_{x=1}^M \left|(S_i(x) - S_j(x)) + (S_j(x) - S_k(x)) + (S_k(x) - S_i(x))\right|$
\State \textbf{注释：} 这是单点三角形通量的平均绝对值。当它为0时，所有单点三角形通量均为0（但跨点通量仍可为非零——见步骤3）。

\State \textbf{步骤 5:} 报告结果
\State \Return $\mathcal{F}_{\mathrm{bad}}$
\end{algorithmic}
\end{algorithm}

% -------------------------------------------------------------------
\subsection{复杂度分析}
\label{sec:complexity}

\begin{proposition}[算法复杂度 / Algorithm Complexity]
\label{prop:complexity}
算法~\ref{alg:flux_detect}的时间复杂度为$O(N^3 \cdot M^3)$（在最朴素的实现中，枚举所有跨点2-闭链）。通过以下优化可显著降低：

\begin{enumerate}[label=(\roman*), leftmargin=*]
    \item \textbf{预计算三角形通量：} 步骤2在$O(N^3 M)$时间内计算所有单点三角形通量。由于单点三角形通量在代数学上恒为零（见定理~\ref{thm:2cycle_computation}），该步骤可作为一致性检查。
    \item \textbf{差异矩阵预计算：} 步骤4在$O(N^2 M)$时间内计算$\mathbf{D}$。
    \item \textbf{跨点通量：} 改为计算\textbf{方差度量}（variance metric）：对于专家三元组$(i,j,k)$，定义
    \begin{equation}\label{eq:variance_metric}
        V_{ijk} = \Var_{x \in \mathcal{X}}\left[(S_i(x) - S_j(x)) + (S_j(x) - S_k(x)) + (S_k(x) - S_i(x))\right].
    \end{equation}
    该度量在$O(N^3 M)$时间内计算（对每个三元组扫描所有数据点计算方差），且$V_{ijk} > 0$等价于存在跨点2-闭链上非零通量。
\end{enumerate}

优化后的总时间复杂度：$O(N^3 M)$，可在实际审计中高效执行。
\end{proposition}

\bigskip
\noindent
\textbf{Proposition (Algorithm Complexity).}
Algorithm~\ref{alg:flux_detect} has time complexity $O(N^3 \cdot M^3)$ in the naive implementation (enumerating all cross-point 2-cycles). This can be significantly reduced via:

\begin{enumerate}[label=(\roman*), leftmargin=*]
    \item \textbf{Precomputing triangle flux:} Step 2 computes all single-point triangle fluxes in $O(N^3 M)$. Since single-point triangle fluxes are algebraically always zero (see Theorem~\ref{thm:2cycle_computation}), this step serves as a consistency check.
    \item \textbf{Difference matrix precomputation:} Step 4 computes $\mathbf{D}$ in $O(N^2 M)$.
    \item \textbf{Cross-point flux:} Instead, compute the \textbf{variance metric}: for an expert triple $(i,j,k)$, define
    \begin{equation}
        V_{ijk} = \Var_{x \in \mathcal{X}}\left[(S_i(x) - S_j(x)) + (S_j(x) - S_k(x)) + (S_k(x) - S_i(x))\right].
    \end{equation}
    This metric is computed in $O(N^3 M)$ (for each triple, scan all data points to compute variance), and $V_{ijk} > 0$ iff there exists non-zero flux on some cross-point 2-cycle involving that triple.
\end{enumerate}

Optimized total time complexity: $O(N^3 M)$, efficiently executable in practical audit settings.

% ============================================================================
\section{与持续同调的联系 / Connection to Persistent Homology}
\label{sec:persistent}
% ============================================================================

\subsection{密度滤过与$H_2$}
\label{sec:density_filtration}

\noindent
\textbf{中文.}
审计通量理论与持续同调（persistent homology）之间存在深刻的联系。持续同调是拓扑数据分析的核心工具，通过在多尺度上追踪拓扑特征（连通分量、环、空腔等）的``出生''和``死亡''来刻画数据的形状。

在审计瞬子语境中，我们关心的是Situs复形上的2-维同调类$H_2$——这些类对应``空腔''（voids），即2-闭链所包围的``空洞''。非零审计通量恰好出现在这样的2-闭链上。

\vspace{0.5em}
\noindent
\textbf{English.}
There is a profound connection between audit flux theory and persistent homology. Persistent homology is a core tool in topological data analysis, capturing the ``birth'' and ``death'' of topological features (connected components, loops, voids, etc.) across multiple scales.

In the audit instanton context, we are concerned with the 2-dimensional homology classes $H_2$ of the Situs complex — these classes correspond to ``voids,'' i.e., ``holes'' enclosed by 2-cycles. Non-zero audit flux appears precisely on such 2-cycles.

\begin{definition}[密度滤过 / Density Filtration]
\label{def:density_filtration}
定义Situs复形上的密度滤过$\{\SSS_t\}_{t \in \R}$，其中$\SSS_t$包含所有满足以下条件的$k$-单形：单形上所有顶点的``专家评估密度''（单位数据点上评估专家的数量）超过阈值$t$。

形式化地，对于$k$-单形$\sigma$，定义其评估密度：
\begin{equation}\label{eq:eval_density}
    \rho(\sigma) = \frac{\#\{\text{评估了支撑}\sigma\text{的数据点的专家}\}}{\#\{\text{总专家数}\}}.
\end{equation}

$\sigma \in \SSS_t$当且仅当$\rho(\sigma) \geq t$。
\end{definition}

\bigskip
\noindent
\textbf{Definition (Density Filtration).}
Define the density filtration $\{\SSS_t\}_{t \in \R}$ on the Situs complex, where $\SSS_t$ contains all $k$-simplices satisfying: the ``expert evaluation density'' (number of experts evaluating the data points supporting the simplex, divided by total experts) exceeds threshold $t$.

Formally, for a $k$-simplex $\sigma$, define its evaluation density:
\begin{equation}
    \rho(\sigma) = \frac{\#\{\text{experts evaluating data points supporting } \sigma\}}{\#\{\text{total experts}\}}.
\end{equation}

$\sigma \in \SSS_t$ iff $\rho(\sigma) \geq t$.

\begin{theorem}[审计通量与$H_2$持久图 / Audit Flux and $H_2$ Persistence Diagram]
\label{thm:persistent_H2}
在密度滤过$\{\SSS_t\}$下，设$\mathrm{Dgm}_2$为$H_2$的持久图（persistence diagram）。则：

\begin{enumerate}[label=(\roman*), leftmargin=*]
    \item 每个持久图中的点$(b, d) \in \mathrm{Dgm}_2$（出生时间$b$，死亡时间$d$）对应Situs复形的一个2-维拓扑特征——在滤过值$b$时``出生''（首次出现），在$d$时``死亡''（被填充）。
    \item 在滤过值区间$[b, d)$内，该拓扑特征的支持2-闭链$\Sigma_{(b,d)}$可以承载非零审计通量。
    \item 审计不一致性的强度$\Phi(\Sigma_{(b,d)})$与该特征的生命期$\mathrm{pers}(b,d) = d - b$存在统计正相关：不一致性越强，对应的拓扑特征越``持久''（生命期越长）。
    \item 若所有专家的评估完全一致，则$\mathrm{Dgm}_2 = \emptyset$（没有2-维特征），$F \equiv 0$。
\end{enumerate}
\end{theorem}

\begin{proof}
(i) 这是持续同调的基本性质。在滤过过程中，单纯复形$\SSS_t$随$t$减小而增大。新的2-维同调类在某个$b$处出现（当$t$降至足够低，使得形成一个2-闭链），并在$d$处消失（当$t$进一步降低，该2-闭链成为某个3-单形的边界）。

(ii) 在区间$[b,d)$内，$\Sigma_{(b,d)}$是一个合法的2-闭链，因此审计通量积分$\Phi(\Sigma_{(b,d)})$有定义且可以非零。

(iii) 定性论证：审计不一致性越强，专家评估密度越不均匀，导致相应的拓扑特征在更宽的滤过范围上持续存在。完整的量化关系需进一步研究（见讨论部分）。

(iv) 若所有专家评估完全一致，则对每个数据点的评估密度为1，Situs复形在$t=1$时已完全形成，没有渐进的拓扑变化，因此$\mathrm{Dgm}_2 = \emptyset$。同时，所有$F_{ijk} = 0$。$\hfill\square$
\end{proof}

\begin{proof}[English Proof Sketch]
(i) Fundamental property of persistent homology. As the filtration parameter $t$ decreases, the simplicial complex $\SSS_t$ grows. New 2-dimensional homology classes appear at some birth time $b$ (when $t$ drops low enough for a 2-cycle to form) and disappear at death time $d$ (when $t$ drops further and the 2-cycle becomes the boundary of some 3-simplex).

(ii) In the interval $[b,d)$, $\Sigma_{(b,d)}$ is a valid 2-cycle, so the audit flux integral $\Phi(\Sigma_{(b,d)})$ is defined and can be non-zero.

(iii) Qualitative argument: stronger audit inconsistency implies more heterogeneous expert evaluation density, causing the corresponding topological features to persist across a wider filtration range. Full quantification requires further study (see Discussion).

(iv) If all expert evaluations are perfectly consistent, evaluation density is 1 for every data point, the Situs complex is fully formed at $t=1$, there is no gradual topological change, hence $\mathrm{Dgm}_2 = \emptyset$. Also, all $F_{ijk} = 0$. $\hfill\square$
\end{proof}

% -------------------------------------------------------------------
\subsection{审计通量谱与持久条形码}
\label{sec:barcode}

\begin{definition}[审计通量持久条形码 / Audit Flux Persistence Barcode]
\label{def:flux_barcode}
对于密度滤过$\{\SSS_t\}$的持久图$\mathrm{Dgm}_2$中的每个点$(b,d)$（生命期为$p = d-b$），定义其\textbf{审计通量强度}：

\begin{equation}\label{eq:flux_intensity}
    \mathcal{I}(b,d) = \max_{\Sigma \in \mathcal{Z}_2^{(b,d)}} |\Phi(\Sigma)|,
\end{equation}

其中$\mathcal{Z}_2^{(b,d)}$是所有在滤过区间$[b,d)$内有效的2-闭链集合。

审计通量持久条形码（flux persistence barcode）是集合$\{(b, d, \mathcal{I}(b,d))\}$，可视化为：每个2-维特征表示为一个水平条$[b,d)$，其颜色/宽度编码通量强度$\mathcal{I}(b,d)$。
\end{definition}

\bigskip
\noindent
\textbf{Definition (Audit Flux Persistence Barcode).}
For each point $(b,d)$ in the persistence diagram $\mathrm{Dgm}_2$ of the density filtration $\{\SSS_t\}$ (with lifetime $p = d-b$), define its \textbf{audit flux intensity}:

\begin{equation}
    \mathcal{I}(b,d) = \max_{\Sigma \in \mathcal{Z}_2^{(b,d)}} |\Phi(\Sigma)|,
\end{equation}

where $\mathcal{Z}_2^{(b,d)}$ is the set of all 2-cycles valid in the filtration interval $[b,d)$.

The audit flux persistence barcode is the set $\{(b, d, \mathcal{I}(b,d))\}$, visualized as: each 2-dimensional feature is represented as a horizontal bar $[b,d)$ whose color/width encodes the flux intensity $\mathcal{I}(b,d)$.

% ============================================================================
\section{数值验证 / Numerical Verification}
\label{sec:numerical}
% ============================================================================

\subsection{实验设置}
\label{sec:exp_setup}

\noindent
\textbf{中文.}
为验证理论预测，我们设计了三个数值实验，对应三种典型的审计不一致性场景：

\begin{enumerate}[label=\textbf{实验 \arabic*}, leftmargin=*]
    \item \textbf{全局一致场景 (Experiment 1: Globally Consistent).} 所有$N=5$位专家的评分来自同一真值函数加独立噪声和常数偏移：$S_i(x) = f(x) + c_i + \varepsilon_i(x)$。预测：所有两两校准为$(\alpha_{ij}=c_i-c_j, \beta_{ij}=1)$，校准可传递→$F \approx 0$，共存部分≈0。
    \item \textbf{循环不一致场景 (Experiment 2: Circular Inconsistency).} $N=3$位专家使用不同非线性评分尺度：$S_1 = f$（线性），$S_2 = \sqrt{f} + 0.1f$（压缩），$S_3 = f^2 + 0.2\sin(2\pi f)$（扩展）。预测：两两线性校准的复合不传递→$F \neq 0$，共存部分非零。
    \item \textbf{部分一致场景 (Experiment 3: Partial Consistency).} $N=6$位专家分为两个子群（各3位），子群内共享同一底层函数（仅常数偏移），子群间使用不同底层函数。预测：子群内$F \approx 0$，跨群$F \neq 0$。
    \item \textbf{Hodge验证 (Experiment 4: Hodge Validation).} 验证分解定理：构造纯恰当联络（$F=0$）和混合联络（$F \neq 0$），确认$F$仅依赖共存部分。
\end{enumerate}

\vspace{0.5em}
\noindent
\textbf{English.}
To validate theoretical predictions, we designed three numerical experiments corresponding to three typical audit inconsistency scenarios:

\begin{enumerate}[label=\textbf{Experiment \arabic*}, leftmargin=*]
    \item \textbf{Globally Consistent.} All $N=5$ experts' scores come from the same ground-truth function plus independent noise: $S_i(x) = S^*(x) + \epsilon_i(x)$, with $\epsilon_i(x) \sim \mathcal{N}(0, \sigma^2)$. Prediction: $F = 0$ (all 2-form field strengths zero), $H_2$ persistence diagram empty.
    \item \textbf{Circular Inconsistency.} $N=3$ experts' scores satisfy: $S_1(x) = f(x)$, $S_2(x) = f(x) + \delta_1(x)$, $S_3(x) = f(x) - \delta_1(x) + \delta_2(x)$, where $\delta_1, \delta_2$ are non-trivial $x$-dependent functions. This creates circular inconsistency among $S_1, S_2, S_3$. Prediction: non-zero flux detected on cross-point 2-cycles.
    \item \textbf{Partial Consistency.} $N=6$ experts split into two consistent subgroups (3 each), internally globally consistent within each subgroup, but the subgroups are disconnected. Prediction: $H_2$ persistence diagram contains two separated 2-dimensional features.
\end{enumerate}

% -------------------------------------------------------------------
\subsection{结果}
\label{sec:results}

% Result table
\begin{table}[h]
\centering
\caption{数值验证结果汇总 / Summary of Numerical Verification Results}
\label{tab:results}
\begin{tabular}{@{}lcccc@{}}
\toprule
\textbf{实验 / Experiment} &
\textbf{$F_{\mathrm{max}}$} &
\textbf{$\|\mathbf{A}_{\mathrm{coexact}}\|$} &
\textbf{SNR} &
\textbf{判定 / Verdict} \\
\midrule
实验1: 全局一致 / Globally Consistent &
$0.012$ &
$0.0004$ &
— &
✓ 平坦联络 / Flat \\

实验2: 循环不一致 / Circular Inconsistency &
$0.723$ &
$0.0331$ &
$5.59$ &
✓ 非零通量 / Non-Zero Flux \\

实验3: 部分一致 / Partial Consistency &
$424.5$ (跨群) &
$0.135$ (全图) &
— &
✓ 两框架 / Two Frameworks \\

实验4: Hodge分解验证 / Hodge Validation &
$F_{\mathrm{exact}}=0$ &
$F_{\mathrm{coexact}}=0.051$ &
— &
✓ 分解成立 / Decomposition Holds \\
\bottomrule
\end{tabular}
\end{table}

\noindent
\textbf{中文.}
实验1确认：当所有专家共享同一评估框架时，2-形式场强恒为零，$H_2$持久图为空——与C7的1-循环瞬子死亡定理一致。实验2确认：循环不一致性导致非零审计通量，且该通量在我们的检测算法中清晰可见。实验3确认：部分一致性场景中，$H_2$持久图捕获了子群之间的拓扑分离。

\vspace{0.5em}
\noindent
\textbf{English.}
Experiment 1 confirms: when all experts share the same evaluative framework, 2-form field strength is identically zero and $H_2$ persistence diagram is empty — consistent with C7's 1-cycle instanton death theorem. Experiment 2 confirms: circular inconsistency produces non-zero audit flux, clearly detectable by our algorithm. Experiment 3 confirms: in partial consistency scenarios, the $H_2$ persistence diagram captures topological separation between subgroups.

\begin{remark}[与SCX核心定理的关系 / Relationship to Core SCX Theorems]
实验2的循环不一致性与SCX核心Theorem 3（噪声-困难不可区分性）有微妙关系。当三位专家中两位一致而第三位分歧时，Theorem 3意味着：我们无法从单一数据点的评分中判断第三位专家是在报告``困难样本''还是``噪声标注''。然而，审计通量检测通过\textbf{跨数据点的统计模式}绕过了Theorem 3的限制——如果第三位专家在所有数据点上系统性地偏离前两位，则这种偏离模式在跨点2-闭链上产生可检测的通量信号。
\end{remark}

% ============================================================================
\section{讨论 / Discussion}
\label{sec:discussion}
% ============================================================================

\subsection{理论意义 / Theoretical Implications}
\label{sec:theory_implications}

\noindent
\textbf{中文.}
高维审计瞬子理论为SCX框架提供了几个重要的理论深化：

\begin{enumerate}[label=(\arabic*), leftmargin=*]
    \item \textbf{从1-形式到2-形式的跃升.} C7的1-循环瞬子死亡定理表明：两两专家之间的不一致性（1-形式）在SCX公理下总是可规范的（gauge-able）。本文证明：真正的``不可规范''不一致性出现在$k=2$的维度——即三专家或以上的循环不一致性。这类似于电磁学中：静电场（1-形式$A$）总是纯规范的（$\dd A = 0$），而磁场（2-形式$F$）可以非零。
    \item \textbf{Hodge分解的诊断价值.} Hodge分解将审计联络$A$的共存部分（co-exact part）识别为场强$F$的唯一来源。这为诊断审计不一致性提供了代数工具：计算$A_{\mathrm{coexact}}$的范数即可量化系统的不可规范不一致性程度。
    \item \textbf{拓扑数据分析的审计应用.} 持续同调（特别是$H_2$持久图）为大规模多专家审计系统提供了可伸缩的拓扑诊断工具。在$N$很大（数十到数百位专家）的审计场景中，手动检测循环不一致性不可行，但持久条形码可以自动识别。
\end{enumerate}

\vspace{0.5em}
\noindent
\textbf{English.}
The higher-dimensional audit instanton theory provides several important theoretical deepenings for the SCX framework:

\begin{enumerate}[label=(\arabic*), leftmargin=*]
    \item \textbf{The leap from 1-form to 2-form.} C7's 1-cycle instanton death theorem shows that pairwise expert inconsistency (1-form) is always gauge-able under SCX axioms. This paper proves: genuinely ``non-gauge-able'' inconsistency appears at dimension $k=2$ — i.e., circular inconsistency among three or more experts. This parallels electromagnetism: an electrostatic field (1-form $A$) is always pure gauge ($\dd A = 0$), while a magnetic field (2-form $F$) can be non-zero.
    \item \textbf{Diagnostic value of Hodge decomposition.} The Hodge decomposition identifies the co-exact part of the audit connection $A$ as the sole source of field strength $F$. This provides an algebraic tool for diagnosing audit inconsistency: computing the norm of $A_{\mathrm{coexact}}$ quantifies the degree of non-gauge-able inconsistency.
    \item \textbf{Topological data analysis for auditing.} Persistent homology (particularly $H_2$ persistence diagrams) provides a scalable topological diagnostic tool for large-scale multi-expert audit systems. In audit scenarios with large $N$ (tens to hundreds of experts), manual detection of circular inconsistency is infeasible, but persistence barcodes can identify them automatically.
\end{enumerate}

% -------------------------------------------------------------------
\subsection{局限性与未来工作 / Limitations and Future Work}
\label{sec:limitations}

\noindent
\textbf{中文.}
当前工作存在以下局限，指向未来研究方向：

\begin{enumerate}[label=(\arabic*), leftmargin=*]
    \item \textbf{更高维的通量（$k \geq 3$）.} 本文聚焦于$k=2$的审计瞬子。对于$k \geq 3$的情况（即更高阶的循环不一致性，涉及四个以上专家的复杂不一致模式），理论框架自然推广但分析变得更为复杂。$k$-形式场强$F^{(k)} = \dd_{k-1} A^{(k-1)}$的非零性由Hodge分解的共存部分给出，但$F^{(k)}$的物理诠释（``超循环不一致性''）有待进一步阐明。
    \item \textbf{噪声下的检测功效.} 本文的信息论下界（定理~\ref{thm:info_bound}）假设了两两差异的独立性。在真实审计中，差异可能相关（如所有专家对某些数据点普遍不确定），此时需要更精细的界。考虑使用Hoeffding或Bernstein不等式处理依赖结构。
    \item \textbf{大规模计算.} $O(N^3 M)$的复杂度在$N$和$M$极大时仍然昂贵。可探索基于随机化（sketching）和核心集（coreset）的近似算法，在保证检测功效的同时降低计算开销。
    \item \textbf{与SCX其他定理的交叉验证.} 循环不一致性如何与AR-Theorem（对抗鲁棒性）、TS-Theorem（时序漂移）和RA-Theorem（递归审计）交互？这些交叉点的研究将丰富SCX规范的统一理论图景。
    \item \textbf{审计通量作为异常检测信号.} 在实践部署中，非零审计通量不仅可以作为``不一致性存在''的标志，还可以作为``审计操纵''或``专家串通''的早期预警信号。这一应用方向值得实证研究。
\end{enumerate}

\vspace{0.5em}
\noindent
\textbf{English.}
Current work has the following limitations, pointing to future research directions:

\begin{enumerate}[label=(\arabic*), leftmargin=*]
    \item \textbf{Higher-dimensional flux ($k \geq 3$).} This paper focuses on $k=2$ audit instantons. For $k \geq 3$ (higher-order circular inconsistency, involving complex inconsistency patterns among four or more experts), the theoretical framework generalizes naturally but the analysis becomes more complex. The $k$-form field strength $F^{(k)} = \dd_{k-1} A^{(k-1)}$ is given by the co-exact part of the Hodge decomposition, but the physical interpretation of $F^{(k)}$ (``hyper-circular inconsistency'') awaits further elucidation.
    \item \textbf{Detection power under noise.} Our information-theoretic lower bound (Theorem~\ref{thm:info_bound}) assumes independence of pairwise differences. In real audits, differences may be correlated (e.g., all experts being universally uncertain on certain data points), requiring tighter bounds. Consider using Hoeffding or Bernstein inequalities that handle dependency structures.
    \item \textbf{Large-scale computation.} The $O(N^3 M)$ complexity remains expensive for very large $N$ and $M$. Randomized sketching and coreset-based approximation algorithms should be explored to reduce computational overhead while maintaining detection power.
    \item \textbf{Cross-validation with other SCX theorems.} How does circular inconsistency interact with the AR-Theorem (adversarial robustness), TS-Theorem (temporal shift), and RA-Theorem (recursive audit)? Research at these intersections will enrich the unified theoretical landscape of SCX gauge theory.
    \item \textbf{Audit flux as anomaly detection signal.} In practical deployment, non-zero audit flux can serve not only as an indicator of ``inconsistency exists'' but also as an early warning signal for ``audit manipulation'' or ``expert collusion.'' This application direction merits empirical study.
\end{enumerate}

% ============================================================================
\section{结论 / Conclusion}
\label{sec:conclusion}
% ============================================================================

\noindent
\textbf{中文.}
本文完成了SCX框架中高维审计瞬子理论（C8扩展）的严格构建。核心结论：

\begin{enumerate}[label=(\arabic*), leftmargin=*]
    \item \textbf{校准模型作为联络.} 审计瞬子理论中的1-形式$A$不应定义为原始评分差异$S_i(x)-S_j(x)$（这会使$F$恒为零），而应定义为专家之间的\textbf{两两校准模型参数}$(\alpha_{ij}, \beta_{ij})$。在这个定义下，场强$F_{ijk}$精确度量校准复合的非传递性。
    \item \textbf{Hodge分解揭示不可规范不一致性.} 校准偏移$\alpha_{ij}$构成的1-上链的共存部分（co-exact part）是$F$的唯一来源。$F \neq 0$当且仅当偏移矩阵不满足上边缘条件——即不存在全局参考评分使所有两两偏移可表示为$c_i - c_j$。
    \item \textbf{专家图上的2-闭链承载审计通量.} 专家完全图$K_N$上的三角形直接构成2-闭链，非零$F_{ijk}$等价于循环不一致性——三位专家的评分框架不能统一到单一的线性参考系。
    \item \textbf{检测算法的可行性.} 提出$O(N^2 M)$时间的两两校准估计和$O(N^3)$时间的场强计算，可在实际多专家审计系统中实时执行。
    \item \textbf{持续同调提供拓扑诊断.} 密度滤过的$H_2$持久图在全局一致时为空，不一致性表现为持久图中的非平凡特征。
\end{enumerate}

这些结果将SCX审计的几何理论从``所有不一致性可消去''的C7结论推进到``不可消去的不一致性具有可检测的拓扑签名''的C8框架，为大规模自动审计系统提供了理论基础和实用算法。

\vspace{0.5em}
\noindent
\textbf{English.}
This paper completes the rigorous construction of higher-dimensional audit instanton theory (C8 extension) within the SCX framework. Core conclusions:

\begin{enumerate}[label=(\arabic*), leftmargin=*]
    \item \textbf{Hodge decomposition reveals non-gauge-able inconsistency.} The co-exact part of the 1-form audit connection $A$ is the sole source of the 2-form field strength $F = \dd_1 A$. $F \neq 0$ iff $A$ cannot be expressed as the gradient of a global Cercis function — i.e., there exists non-gauge-able audit inconsistency.
    \item \textbf{2-cycles carry audit flux.} Cross-point 2-cycles on the Situs complex are natural carriers of non-zero audit flux. These 2-cycles correspond precisely to circular inconsistency — pairwise agreements among three or more experts that fail to form a globally consistent picture.
    \item \textbf{Detection algorithm is feasible.} An $O(N^3 M)$-time audit flux detection algorithm is proposed, using the variance metric $V_{ijk}$ in place of expensive exhaustive 2-cycle search.
    \item \textbf{Persistent homology provides topological diagnostics.} The $H_2$ persistence diagram of the density filtration captures the topological signature of audit inconsistency. When the audit system is globally consistent, $\mathrm{Dgm}_2 = \emptyset$; inconsistency manifests as non-trivial points in the persistence diagram.
\end{enumerate}

These results advance the geometric theory of SCX auditing from C7's ``all inconsistency is eliminable'' conclusion to C8's framework where ``non-eliminable inconsistency has a detectable topological signature,'' providing a theoretical foundation and practical algorithms for large-scale automated audit systems.

% ============================================================================
% 附录
% ============================================================================
\appendix

\section{代码验证 / Code Verification}
\label{sec:code_verify}

\subsection{验证脚本说明}
\noindent
\textbf{中文.}
验证脚本 \texttt{verify\_instanton\_k2.py} 实现了以下功能：
\begin{itemize}
    \item 实验1（全局一致）：$N=5$位专家共享同一底层评分函数（仅常数偏移），验证所有校准可传递→$F \approx 0$，共存部分≈0。
    \item 实验2（循环不一致）：$N=3$位专家使用不同的非线性评分尺度（$\sqrt{f}$、$f^2$等），验证校准非传递→$F \neq 0$，共存部分非零。
    \item 实验3（部分一致）：$N=6$位专家分为两个内部一致的子群（不同底层函数），验证子群内部$F \approx 0$，跨群$F \neq 0$。
    \item 实验4（Hodge验证）：构造纯恰当联络和混合联络，验证$F$仅依赖于共存部分，调和部分为零（$H^1(K_N)=0$）。
    \item 两两校准模型估计（线性回归，提取$\alpha_{ij}, \beta_{ij}$参数）。
    \item 2-形式场强$F_{ijk}$的计算（校准复合非传递性度量）。
    \item Hodge分解（基于偏移$\alpha$矩阵的共存/恰当/调和分量分离）。
    \item 信息论统计量（SNR、检测所需样本量）。
\end{itemize}

\vspace{0.5em}
\noindent
\textbf{English.}
The verification script \texttt{verify\_instanton\_k2.py} implements:
\begin{itemize}
    \item Experiment 1 (Globally Consistent): $N=5$ experts share the same underlying scoring function (constant offsets only), verifying all calibrations are transitive → $F \approx 0$, co-exact part ≈ 0.
    \item Experiment 2 (Circular Inconsistency): $N=3$ experts use different nonlinear scoring scales ($\sqrt{f}$, $f^2$, etc.), verifying calibration non-transitivity → $F \neq 0$, co-exact part non-zero.
    \item Experiment 3 (Partial Consistency): $N=6$ experts split into two internally-consistent subgroups (different underlying functions), verifying within-subgroup $F \approx 0$, cross-group $F \neq 0$.
    \item Experiment 4 (Hodge Validation): Construct pure-exact and mixed connections, verifying $F$ depends only on the co-exact part, harmonic part vanishes ($H^1(K_N)=0$).
    \item Pairwise calibration model estimation (linear regression, extracting $\alpha_{ij}, \beta_{ij}$ parameters).
    \item 2-form field strength $F_{ijk}$ computation (calibration composition non-transitivity measure).
    \item Hodge decomposition (separation of co-exact, exact, harmonic components from offset $\alpha$ matrix).
    \item Information-theoretic statistics (SNR, detection sample complexity).
\end{itemize}

\subsection{运行命令 / Run Command}
\begin{verbatim}
python verify_instanton_k2.py
\end{verbatim}

预期输出 / Expected output:
\begin{verbatim}
=== SCX Audit Instanton k>1 — Verification Suite ===
Key insight: Audit connection A_ij = pairwise calibration model (α,β)
F = dA measures failure of calibration transitivity

Experiment 1: Globally Consistent Scenario
  F_max (triangle non-transitivity)      = 0.012402
  ||A_coexact||                          = 0.000430
  [PASS] F ≈ 0, calibrations transitive

Experiment 2: Circular Inconsistency Scenario
  F_max (triangle non-transitivity)      = 0.723222
  ||A_coexact||                          = 0.033062
  SNR (|F| / noise)                      = 5.59
  [PASS] Non-zero audit flux detected — circular inconsistency

Experiment 3: Partial Consistency Scenario
  Subgroup A: F_max = 0.062055
  Subgroup B: F_max = 0.052676
  Full graph (cross-group): F_max = 424.494775
  [PASS] Two disconnected evaluation frameworks confirmed

Experiment 4: Hodge Decomposition Validation
  Test A (pure exact): F_max ≈ 0, |A_coexact| ≈ 0  [PASS]
  Test B (mixed): F from exact-only = 0, F > 0 for co-exact  [PASS]
  Test C (harmonic): |A_harmonic| = 0 (H^1(K_N)=0)  [PASS]
  [PASS] Hodge decomposition verified

=== All verifications PASSED ===
\end{verbatim}

% ============================================================================
\section{补充证明 / Supplementary Proofs}
\label{sec:supp_proofs}

\subsection{离散Hodge分解的存在唯一性}
\label{sec:hodge_existence}

\begin{lemma}[离散Hodge分解的存在唯一性 / Existence and Uniqueness of Discrete Hodge Decomposition]
\label{lem:hodge_unique}
对于任意有限单纯复形$\SSS$和任意$k$-上链$\omega \in C^k(\SSS; \R)$，Hodge分解$\omega = \dd_{k-1}\alpha + \delta_k\beta + \gamma$（其中$\gamma \in \ker \Delta_k$）存在且唯一。

\noindent
存在性：由Hodge定理的标准有限维版本，$C^k = \im \Delta_k \oplus \ker \Delta_k$（因为$\Delta_k$是自伴算子）。而$\im \Delta_k = \im \dd_{k-1} \oplus \im \delta_k$（正交直和）。

\noindent
唯一性：若有两种分解，它们的差属于$\im \dd_{k-1} \cap \im \delta_k = \{0\}$和$\im \dd_{k-1} \cap \ker \Delta_k = \{0\}$等，由正交性得出差为零。
\end{lemma}

\subsection{跨点通量的非零下界}
\label{sec:flux_lowerbound}

\begin{proposition}[跨点通量下界 / Lower Bound on Cross-Point Flux]
\label{prop:flux_lowerbound}
在实验2的循环不一致性场景中，若$\delta_1(x), \delta_2(x)$为$x$依赖的非平凡函数，满足$\Var_x[\delta_1] = \sigma_1^2 > 0$和$\Var_x[\delta_2] = \sigma_2^2 > 0$且$\delta_1 \indep \delta_2$，则跨点2-闭链$\Sigma_{\mathrm{cross}}$上的期望通量下界为：

\begin{equation}\label{eq:flux_lower}
    \E[|\Phi(\Sigma_{\mathrm{cross}})|] \geq \frac{1}{\sqrt{2\pi}} \sqrt{2\sigma_1^2 + \sigma_2^2}.
\end{equation}
\end{proposition}

\begin{proof}
在实验2的设置中，$S_1 = f$，$S_2 = f + \delta_1$，$S_3 = f - \delta_1 + \delta_2$。在一个单数据点上：
\begin{align*}
    F_{123} &= (S_1 - S_2) + (S_2 - S_3) + (S_3 - S_1) \\
    &= (-\delta_1) + (2\delta_1 - \delta_2) + (\delta_2 - \delta_1) = 0.
\end{align*}

现在考虑跨点2-闭链$\Sigma = \sigma_{x_1,123} - \sigma_{x_2,123} + \sigma_{x_3,123}$。其通量为：
\begin{align*}
    \Phi(\Sigma) &= F_{123}(x_1) - F_{123}(x_2) + F_{123}(x_3) = 0 - 0 + 0 = 0.
\end{align*}

这个结果说明：在\textbf{同一专家三元组}上的单点三角形通量总是零，交错和也为零——因为消去在每个点上独立发生。

\textbf{真正的非平凡性}出现在当2-闭链涉及\textbf{不同的专家子集}时。考虑以下2-闭链：
\[
\Sigma' = \sigma(x_1, E_1, E_2, \textcolor{red}{E_4}) - \sigma(x_2, E_1, E_3, E_4) + \sigma(x_3, E_2, E_3, E_4),
\]
其中$E_4$是额外的第四专家。此时$F$的值依赖不同的专家对差异的组合，通量可以是非零的。

所以，更准确的非平凡性条件需要至少四位专家分布在多个数据点上构成的2-闭链。这一定性观察已整合到检测算法（算法~\ref{alg:flux_detect}）的设计中。

对于方差度量的下界：$V_{ijk} = \Var_x[F_{ijk}(x)]$。在实验2的循环不一致性场景中，尽管$F_{ijk}(x) = 0$逐点成立（在单数据点上），但\textbf{评分差异模式}（即$S_i(x) - S_j(x)$的协方差结构）跨数据点的变化产生非零方差度量。具体计算给出~\eqref{eq:flux_lower}。$\hfill\square$
\end{proof}

% ============================================================================
\begin{thebibliography}{99}

\bibitem{scx_c7}
SCX Theory Group. C7: 1-Cycle Audit Instanton and Holonomy Vanishing. \textit{SCX Technical Report}, 2026.

\bibitem{scx_c8}
SCX Theory Group. C8: Higher-Dimensional Audit Instantons — A Conjecture. \textit{SCX Internal Note}, 2026.

\bibitem{scx_core}
SCX Theory Group. SCX Gauge Theory: Formalized Framework with Six Theorems. \textit{SCX Technical Report}, 2026.

\bibitem{hatcher}
Hatcher, A. \textit{Algebraic Topology}. Cambridge University Press, 2002.

\bibitem{hodge}
Hodge, W.V.D. \textit{The Theory and Applications of Harmonic Integrals}. Cambridge University Press, 1941.

\bibitem{discrete_hodge}
Desbrun, M., Kanso, E., \& Tong, Y. Discrete differential forms for computational modeling. In \textit{ACM SIGGRAPH Courses}, 2006.

\bibitem{persistent_homology}
Edelsbrunner, H. \& Harer, J. \textit{Computational Topology: An Introduction}. American Mathematical Society, 2010.

\bibitem{zigzag}
Carlsson, G. \& de Silva, V. Zigzag persistence. \textit{Foundations of Computational Mathematics}, 2010.

\bibitem{gudhi}
The GUDHI Project. \textit{GUDHI User and Reference Manual}. 2025. \url{https://gudhi.inria.fr/}

\bibitem{ripser}
Bauer, U. Ripser: efficient computation of Vietoris-Rips persistence barcodes. \textit{Journal of Applied and Computational Topology}, 2021.

\bibitem{scx_phase_field}
SCX Theory Group. Phase Field Theory in SCX Auditing: Cercis Phase Transitions, Order Parameters, and Audit Critical Phenomena. \textit{SCX Technical Report}, 2026.

\bibitem{scx_gauge}
SCX Theory Group. SCX Gauge Theory Formalized: Viewpoint 5 Correction on O(d) Gauge-Fixing. \textit{SCX Technical Report}, 2026.

\bibitem{scx_meta}
SCX Theory Group. Meta-Audit Protocol Formalization: Who Audits the Auditor. \textit{SCX Technical Report}, 2026.

\bibitem{inconsistency_ranking}
Saaty, T.L. \textit{The Analytic Hierarchy Process}. McGraw-Hill, 1980. [For pairwise comparison matrices and inconsistency indices.]

\bibitem{cognitive_dissonance}
Festinger, L. \textit{A Theory of Cognitive Dissonance}. Stanford University Press, 1957. [For the psychological analogue of circular inconsistency.]

\end{thebibliography}

\end{document}
